% GENERATED FILE. Edit canonical lessons, then run npm run generate:books.
% canonical-source-hash: fnv1a64:3269c9aaa360cb2a
% canonical-lessons: TA-C150-uta, TA-C150-campal, TA-C150-patu, TA-C150-camai, TA-C150-vettu

\chapter{Purple, Grey, To sing, To cook, To cut}
\label{ch:ta-purple-grey-to-sing-to-cook-to-cut}
\hlchaptermodality{150}

\begin{chapteropening}
\textbf{By the end of this chapter:} \emph{I can say purple, grey, to sing, to cook and to cut.}

Complete the last lesson of chapter 150: i can say purple, grey, to sing, to cook and to cut.
\end{chapteropening}

\section[ūtā]{\ta{ஊதா} (ūtā) — purple}

\label{lesson:TA-C150-uta}



\begin{quote}
Before the new one: say the Tamil for true, then the Tamil for false.
\end{quote}

\subsection*{You'll want to know: \ta{ஊதா}}
\textbf{\ta{ஊதா}} — \emph{ūtā} — \textquotedblleft{}purple\textquotedblright{}.

\begin{grammarlens}[title={What you've built}]
One more word for this chapter.
\end{grammarlens}

\subsection*{Your turn}
\begin{itemize}
\raggedright
  \item \emph{Say it:} \emph{ūtā}
  \item \emph{Say it:} \emph{ūtā}, once more
  \item \emph{Say it:} say \emph{ūtā}
  \item \emph{Recall:} say the Tamil for true, then the Tamil for false, then say \emph{ūtā} again
\end{itemize}

\begin{tcolorbox}[breakable,colback=teal!4,colframe=teal!35!black,title={Before you move on}]
What does \ta{ஊதா} mean? (\textquotedblleft{}Purple\textquotedblright{}.) Say it once more.
\end{tcolorbox}
\section[cāmpal]{\ta{சாம்பல்} (cāmpal) — grey; ash}

\label{lesson:TA-C150-campal}



\begin{quote}
Before the new one: say the Tamil for false, then the Tamil for purple.

Say the colours: \textbf{\ta{கருப்பு}}, \textbf{\ta{வெள்ளை}}, \textbf{\ta{சிவப்பு}}, \textbf{\ta{நீலம்}}, \textbf{\ta{பச்சை}}, \textbf{\ta{மஞ்சள்}}. Blue is the one many languages named last.
\end{quote}

\subsection*{You'll want to know: \ta{சாம்பல்}}
\textbf{\ta{சாம்பல்}} — \emph{cāmpal} — \textquotedblleft{}grey; ash\textquotedblright{}.

\begin{grammarlens}[title={What you've built}]
One more word for this chapter.
\end{grammarlens}

\subsection*{Your turn}
\begin{itemize}
\raggedright
  \item \emph{Say it:} \emph{cāmpal}
  \item \emph{Say it:} \emph{cāmpal}, once more
  \item \emph{Say it:} say \emph{cāmpal}
  \item \emph{Recall:} say the Tamil for false, then the Tamil for purple, then say \emph{cāmpal} again
\end{itemize}

\begin{tcolorbox}[breakable,colback=teal!4,colframe=teal!35!black,title={Before you move on}]
What does \ta{சாம்பல்} mean? (\textquotedblleft{}Grey; ash\textquotedblright{}.) Say it once more.
\end{tcolorbox}
\section[pāṭu]{\ta{பாடு} (pāṭu) — to sing}

\label{lesson:TA-C150-patu}



\begin{quote}
Before the new one: say the Tamil for purple, then the Tamil for grey; ash.
\end{quote}

\subsection*{You'll want to know: \ta{பாடு}}
\textbf{\ta{பாடு}} — \emph{pāṭu} — \textquotedblleft{}to sing\textquotedblright{}.

\begin{grammarlens}[title={What you've built}]
One more word for this chapter.
\end{grammarlens}

\subsection*{Your turn}
\begin{itemize}
\raggedright
  \item \emph{Say it:} \emph{pāṭu}
  \item \emph{Say it:} \emph{pāṭu}, once more
  \item \emph{Say it:} say \emph{pāṭu}
  \item \emph{Recall:} say the Tamil for purple, then the Tamil for grey; ash, then say \emph{pāṭu} again
\end{itemize}

\begin{tcolorbox}[breakable,colback=teal!4,colframe=teal!35!black,title={Before you move on}]
What does \ta{பாடு} mean? (\textquotedblleft{}To sing\textquotedblright{}.) Say it once more.
\end{tcolorbox}
\section[camai]{\ta{சமை} (camai) — to cook}

\label{lesson:TA-C150-camai}



\begin{quote}
Before the new one: say the Tamil for grey; ash, then the Tamil for to sing.
\end{quote}

\subsection*{You'll want to know: \ta{சமை}}
\textbf{\ta{சமை}} — \emph{camai} — \textquotedblleft{}to cook\textquotedblright{}.

\begin{grammarlens}[title={What you've built}]
One more word for this chapter.
\end{grammarlens}

\subsection*{Your turn}
\begin{itemize}
\raggedright
  \item \emph{Say it:} \emph{camai}
  \item \emph{Say it:} \emph{camai}, once more
  \item \emph{Say it:} say \emph{camai}
  \item \emph{Recall:} say the Tamil for grey; ash, then the Tamil for to sing, then say \emph{camai} again
\end{itemize}

\begin{tcolorbox}[breakable,colback=teal!4,colframe=teal!35!black,title={Before you move on}]
What does \ta{சமை} mean? (\textquotedblleft{}To cook\textquotedblright{}.) Say it once more.
\end{tcolorbox}
\section[veṭṭu]{\ta{வெட்டு} (veṭṭu) — to cut}

\label{lesson:TA-C150-vettu}



\begin{quote}
Before the new one: say the Tamil for to sing, then the Tamil for to cook.
\end{quote}

\subsection*{You'll want to know: \ta{வெட்டு}}
\textbf{\ta{வெட்டு}} — \emph{veṭṭu} — \textquotedblleft{}to cut\textquotedblright{}.

\begin{grammarlens}[title={What you've built}]
One more word for this chapter.
\end{grammarlens}

\subsection*{Your turn}
\begin{itemize}
\raggedright
  \item \emph{Say it:} \emph{veṭṭu}
  \item \emph{Say it:} \emph{veṭṭu}, once more
  \item \emph{Say it:} say \emph{veṭṭu}
  \item \emph{Recall:} say the Tamil for to sing, then the Tamil for to cook, then say \emph{veṭṭu} again
\end{itemize}

\begin{tcolorbox}[breakable,colback=teal!4,colframe=teal!35!black,title={Before you move on}]
What does \ta{வெட்டு} mean? (\textquotedblleft{}To cut\textquotedblright{}.) Say it once more.
\end{tcolorbox}
